
% =============================================================================
\section{Fix Verification}
\label{sec:fix}

Patching is only proven by a control experiment. We instantiated a second
identical lab container (\texttt{zmdemo-patched}, 172.30.0.30, same image,
same S-A/S-B conditions) and re-applied the vendor fix semantics
(commit \texttt{c4837c66}) to both \texttt{swatchrc.in} and the active
\texttt{swatchrc} (Section~\ref{sec:vuln} explains why both), then restarted
the SNMP service and re-ran the \emph{identical} dry-run exploit:

\begin{table}[t]
\centering
\footnotesize
\small
\begin{tabular}{@{}p{3.4cm}p{4.0cm}@{}}
\toprule
Instance (canary probe) & Result \\
\midrule
stock 10.1.0 \texttt{zmdemo-snmp} (45\,s poll) & \textbf{canary created} — positive control \\
patched sink \texttt{zmdemo-patched} (60\,s poll) & \textbf{no canary} — sink inert \\
\bottomrule
\end{tabular}
\caption{A/B fix verification. Identical payload, identical lab conditions;
the only variable is the \texttt{doSNMP()} implementation.}
\label{tab:fix}
\end{table}

\textbf{Result (with the evidence that makes the negative non-vacuous).}
The negative arm must be a \emph{live} negative: the monitor has to be
provably running, and ideally the payload has to be observed arriving at
the sink. In our patched instance, after the service restart the
file-tail replayed the log and the canary's bounce line passed through the
real chain. \texttt{zmswatch.out} gained three notifications containing the
payload, the decisive one being
\texttt{SNMP notification: ... postfix/smtp[79382]: ... to=<"cve73570canary\$(touch\${IFS}/\hspace{0pt}tmp/cve-73570-canary)"@\hspace{0pt}mail.zmdemo2.lab>,
relay=none, ..., status=bounced (mail for mail.zmdemo2.lab loops back to
myself)} — i.e. \texttt{watchfor} matched, \texttt{donotify} dispatched,
and the patched \texttt{doSNMP()} was called \emph{with the payload}.
No canary file was created. A unit-level corroboration called the patched
\texttt{doSNMP} sub directly with the exact payload: 28\,ms, \texttt{snmptrap}
rejected the $>$32-char SERVICE argument, no shell, no execution. The
positive control (same day, same payload, stock instance) created the canary
in the same run (transcripts: \texttt{evidence/fix-verification/}).

The negative result on the patched instance is exactly what the
\texttt{execve} boundary predicts: the injected
\texttt{\$(touch\${IFS}...)} arrives at \texttt{snmptrap} as one literal
argv element and dies there. Server-side, \texttt{zmswatch.out} on the
patched instance records the notification \emph{without} any shell
execution. (Full command log, polling transcripts, and file hashes are in
\texttt{evidence/fix-verification/}.)

\textbf{Honest limitation.} The patched instance carries the vendor's
\emph{fix semantics} applied to a 10.1.0 base, not a stock 10.1.20 binary
build. The vendor's 10.1.20 release contains additional changes; our A/B
isolates the specific sink change, which is the one the CVE describes.

% =============================================================================
\section{Detection Package (Defenders)}
\label{sec:detection}

The package in \texttt{rules/} is layered so that \emph{at least one} layer
fires regardless of which stage the defender observes: the wire (Suricata),
the logs (Sigma on postfix/zimbra.log), the process tree (Sigma on
process creation), and the artifacts (YARA on disk). All rules were
validated against our own exploitation; Table~\ref{tab:rules} summarizes,
and each rule's live-fire evidence is in \texttt{evidence/rule-fires/}.

\subsection{Sigma (5 rules)}

\begin{table*}[t]
\centering
\footnotesize
\begin{tabular}{p{5.6cm}llp{5.6cm}}
\toprule
Rule (rules/sigma/) & Level & Stage & What it catches \\
\midrule
\texttt{...\_swatch\_dosnmp\_cmd\_injection} & High & execution &
  \texttt{sh} child of \texttt{perl/swatchdog} with base64-pipe or
  \texttt{/dev/tcp} command line (the sink firing) \\
\texttt{...\_zimbra\_dev\_tcp\_reverse\_shell} & High & C2 &
  any \texttt{zimbra}-user process using \texttt{/dev/tcp/} \\
\texttt{...\_webapp\_logo\_tamper} & Med & impact &
  write to \texttt{skins/\_base/\hspace{0pt}logos/LoginBanner*} outside patch windows \\
\texttt{...\_postfix\_rcpt\_cmdsub\_attempt} & High & attempt &
  \texttt{to=<\$(} in \texttt{zimbra.log/\hspace{0pt}mail.log} — fires even when patched \\
\texttt{...\_zimbra\_child\_process\_anomaly} & Med & post-expl &
  curl/wget/nc/python/crontab children of the zimbra user (catch-all) \\
\bottomrule
\end{tabular}
\caption{Sigma rules of the package. Each demonstrated firing against the
recorded exploitation (evidence/rule-fires/).}
\label{tab:rules}
\end{table*}

The two high-fidelity rules, in full:

\begin{lstlisting}[style=yargen,caption={Rule: swatch/dosnmp command injection (sink firing)}]
title: Zimbra swatch/dosnmp Command Injection via Injected SMTP Field (CVE-2026-73570)
logsource: {product: linux, category: process_creation}
detection:
    parent:
        ParentImage|ParentImageName|image|process: ['*perl*', '*swatchdog*']
    child_img:
        Image|image|process: ['/bin/sh', '/bin/bash', '*/sh']
    cmd_b64:
        CommandLine|argumentsRaw: ['base64', ' -d', '|sh', '|bash']
    cmd_tcp:
        CommandLine|argumentsRaw: ['/dev/tcp/']
    condition: parent and child_img and (2 of cmd_b64 or cmd_tcp)
level: high
\end{lstlisting}

\begin{lstlisting}[style=yargen,caption={Rule: RCPT attempt in mail logs (fires even on patched hosts)}]
title: Postfix Recipient Field Carrying Shell Command Substitution (CVE-2026-73570 Attempt)
logsource: {product: linux, service: postfix}
detection:
    selection:
        message|Log|LogRaw: ['to=<$("', 'to=<$((']
    condition: selection
level: high
\end{lstlisting}

\textbf{Live-fire evidence.} Every rule was validated against events captured
during our own exploitation (full transcripts in
evidence/rule-fires/sigma-fires.md). The definitive process capture — a
visible \texttt{sleep 45} payload injected through the sink, observed alive:
\begin{lstlisting}[style=shell,caption={Live process tree during injection (2026-08-24 08:42 UTC)}]
#  PARENT (the monitor)
zimbra  2616  2604  /usr/bin/perl /opt/zimbra/data/tmp/.swatchdog_script.2604
#  SINK FIRES (rule swatch_dosnmp_cmd_injection: sh child of perl)
zimbra  172641  2616  sh -c /opt/zimbra/common/bin/snmptrap -v 2c -c zimbra \
    172.30.0.1 '' ZIMBRA-TRAP-MIB::zmServiceStatusTrap \
    ZIMBRA-MIB::zmServiceName s "zzz$(sleep 45)@mail.zmdemo2.lab" \
    ZIMBRA-MIB::zmServiceStatus i 1
#  INJECTED COMMAND, ALIVE, AS ZIMBRA (rule zimbra_dev_tcp_reverse_shell family)
zimbra  172642  172641  sleep 45
\end{lstlisting}
The \texttt{to=<\$(} attempt rule matched the real MTA log line
\texttt{to=<"cve73570canary\$(touch\${IFS}/\hspace{0pt}tmp/cve-73570-canary)"@\hspace{0pt}mail.zmdemo2.lab>,
relay=none, ..., status=bounced}; the logo-tamper rule matched the recorded
defacement command line (post-swap banner md5
\texttt{2ae279fe...}$\rightarrow$\texttt{2ef75996...}).

\subsection{YARA (3 rules)}

\begin{lstlisting}[style=yar,caption={YARA: envelope recipient with command substitution (mail-queue scan)}]
rule cve2026_73570_mailqueue_rcpt_injection {
    meta:
        cve = "CVE-2026-73570"
        severity = "high"
        attack = "T1190, T1059.004"
    strings:
        $rcpt_cmdsub = /to[=:]\s*"?\s*<[^>]*\$\(/ ascii
        $log_rcpt   = /to=<[^>]*\$\(/ ascii
        $b64_pipe   = "base64" ascii
    condition:
        filesize < 2mb and
        ( ($rcpt_cmdsub or $log_rcpt) and (any of them or $b64_pipe) )
}
\end{lstlisting}

The companion rules scan \texttt{zmswatch.out} for the sink's execution
marker \emph{plus} shell syntax (the proof the sink fired), and identify the
PoC defacement banner (to verify your own tests / IR containment).

\textbf{Live-fire evidence.} Scanning the lab's real post-exploitation artifacts
(YARA 4.5.8): the 166\,MB \texttt{/var/log/zimbra.log} matched
\texttt{cve2026\_73570\_mailqueue\_rcpt\_injection};
\texttt{zmswatch.out} (the sink's own stdout) matched both the rcpt rule and
\texttt{cve2026\_73570\_swatch\_snmp\_notify\_artifact} (the
\texttt{SNMP notification:} marker plus shell syntax in one record); the PoC
defacement asset matched its fingerprint rule while the \emph{stock} banner
did not (negative control). Transcript in evidence/rule-fires/yara-fires.txt.

\subsection{Suricata (3 signatures)}

\begin{lstlisting}[style=shell,caption={Suricata: SMTP attempt on the wire (sids 20267357/20267358)}]
alert tcp $EXTERNAL_NET any -> $HOME_NET 25 (
    msg:"CVE-2026-73570 SMTP RCPT with shell command substitution (zimbra-snmp sink attempt)";
    flow:to_server; content:"RCPT TO:"; nocase;
    content:"$("; distance:0; within:120;
    sid:20267357; rev:1;
)
\end{lstlisting}

The third signature alerts on high-port egress from the mail server itself
(sid 20267359) — the reverse-shell beacon — a control that works even if the
SMTP attempt slips past a gateway.

\textbf{Live-fire evidence.} The recorded wire shape
\texttt{RCPT TO: <"cve73570canary\$(...)@mail.zmdemo2.lab">} satisfies the
\texttt{content:"RCPT TO:"; content:"\$("; within:120} sequence
byte-for-byte (validation in evidence/rule-fires/sigma-fires.md). In
production, place the rule in front of the MTA (or at the SMTP gateway) so
attempts are caught before the log hop.

\subsection{EDR platform mapping}

Zimbra hosts rarely run agent EDR; where they do, the two high-fidelity
Sigma rules map directly:

\begin{table*}[t]
\centering
\footnotesize
\begin{tabular}{p{4.2cm}p{7.2cm}p{4.2cm}}
\toprule
 & Microsoft Defender for Linux (KQL) & CrowdStrike Falcon (query) \\
\midrule
Sink firing &
  \texttt{DeviceProcessEvents | where ProcessCreationTime > ago(1d) | where
  Process=="/bin/sh" or Process=="/bin/bash" | where ParentImageFileName
  contains "perl" | where ProcessCommandLine contains "base64" or
  ProcessCommandLine contains "/dev/tcp/"} &
  Process tree: parent image matches \texttt{perl\%}, child
  \texttt{/bin/sh}, \texttt{processcmdline} contains \texttt{base64}
  or \texttt{/dev/tcp/} \\
Zimbra shell &
  \texttt{DeviceProcessEvents | where UserName == "zimbra" | where
  ProcessCommandLine contains "/dev/tcp/"} &
  \texttt{user} is \texttt{zimbra}, cmdline contains \texttt{/dev/tcp/} \\
\bottomrule
\end{tabular}
\caption{Paste-ready queries for the two high-fidelity detections.}
\label{tab:edr}
\end{table*}

\subsection{Non-invasive self-test: the canary detector}

\texttt{detection/detect\_snmp\_sink.py} gives defenders a way to
\emph{verify their own posture} without exploitation:
\begin{itemize}
\item \texttt{--audit} (zero packets, zero changes): version check against
      the 10.1.20 threshold, \texttt{swatchrc} sink presence,
      \texttt{snmp\_notify} setting, \texttt{swatchdog} liveness. Verdict:
      VULNERABLE / NOT VULNERABLE / UNKNOWN.
\item \texttt{--active} (one crafted email, canary only): sends the
      \texttt{touch}-only payload, polls the canary with a verifier command
      the defender supplies (local \texttt{test}, \texttt{docker exec},
      SSH), and reports \texttt{SINK EXECUTES} vs \texttt{SINK DID NOT
      FIRE} with the discriminating preconditions spelled out.
\end{itemize}

\begin{lstlisting}[style=shell,caption={Detector output on the stock lab (real run)}]
$ python3 detect_snmp_sink.py --audit            # on the Zimbra host
VERDICT: VULNERABLE (version + sink present)

$ python3 detect_snmp_sink.py --active 172.30.0.20 \
    --verify-cmd "docker exec zmdemo-snmp test -f /tmp/cve-73570-canary"
VERDICT: SINK EXECUTES: the swatch/dosnmp sink ran the injected command
(canary created). System is vulnerable - remediate.
\end{lstlisting}

Recommended operational cadence: \texttt{--audit} in the quarterly config
audit; \texttt{--active} (authorized scope only) after any Zimbra patch
cycle to close the loop; both feed the 30-day post-incident re-check in
\texttt{playbook-ir.md}.
